能量涨落改变飞行时间，角分布则改变给定探测面处的纵向相位。
二者都可能降低系综干涉，但自由幺正传播本身仍保留单个电子的纯度。
若束流中的其它电子也必须被迹掉，库仑相互作用还可能产生真正的多电子关联。
因此应分别建立这些涨落的模型，不能用一个未经计算的“库仑能散”同时代替它们。

\begin{proposition}[固定探测面上的近轴角平均]
取末态能量 $E_\ell=E_0+\ell\hbar\omega$，
$p_\ell=c^{-1}\sqrt{E_\ell^2-m_{\rm e}^2c^4}$。
假定不同通道拥有同一个横向动量 $p_\perp$，且未读出的两个横向分量是独立
零均值高斯变量，每一分量的方差为 $\sigma_\perp^2$。
近轴展开 $p_{z,\ell}=p_\ell-p_\perp^2/(2p_\ell)+O(p_\perp^4/p_\ell^3)$ 给
\begin{equation}
D_{\ell m}^{(\theta)}(L)=
\left[1+\ii\frac{L\sigma_\perp^2}{\hbar}
\left(\frac1{p_\ell}-\frac1{p_m}\right)\right]^{-1}.
\label{eq:feqo-angular-damping}
\end{equation}
若边带总能宽足够窄，则
\begin{equation}
|D_{\ell m}^{(\theta)}|=
\left[1+(\ell-m)^2\left(\frac{L}{L_\theta}\right)^2\right]^{-1/2},
\qquad L_\theta=\frac{p_0^2v_0}{\omega\sigma_\perp^2}
=\frac{v_0}{\omega\sigma_\theta^2}.
\label{eq:feqo-Ltheta-def}
\end{equation}
最后一式用每一角分量的标准差 $\sigma_\theta=\sigma_\perp/p_0$。
\end{proposition}
\begin{proof}
把 $p_{z,\ell}-p_{z,m}$ 的横向修正乘 $L/\hbar$，得到额外相位
$-\ii a p_\perp^2$，其中 $a=L(1/p_\ell-1/p_m)/(2\hbar)$。
单个高斯分量的平均是 $(1+2\ii a\sigma_\perp^2)^{-1/2}$，
两个独立分量相乘就得到第一式。
再用
$\dd(1/p)/\dd E=-1/(p^2v)$，在 $E_0$ 展开得到模长公式。
四阶近轴项的累计相位仍须满足前面的\eqnrefs{eq:paraxial_phase_control_parameter}；
若横向分布与边带能量相关，就必须对联合分布积分，不能相乘一个通道无关的衰减因子。
\end{proof}

\begin{exbox}[能量涨落与角涨落的长度比较]
取 $\SI{200}{\kilo\electronvolt}$ 电子、$\SI{800}{\nano\metre}$ 光，
每一角分量标准差 $\sigma_\theta=\SI{1}{\milli\radian}$。
由上述模型得 $L_\theta=\SI{0.08853}{\metre}$。
若中心能量标准差为 $\SI{1}{\electronvolt}$，前面的 $L_E=\SI{0.05891}{\metre}$。
在 $L=\SI{0.03}{\metre}$ 处，第一谐波的能量平均因子为 $0.878$，
角平均模长为 $0.947$；二者独立时模长乘积为 $0.831$。
但这两个系综平均不等于每一个电子在自由飞行中不可逆地失去相干。
记录能量或角度并进行条件分析，可以改变所平均的分布。
\end{exbox}

\begin{exbox}[静态库仑涨落必须包含空间积分]
给出一个明确可算的静态模型：在随束团移动的参考系中，
其它粒子的位置构成密度为 $n_*$ 的空间 Poisson 过程，
只计算距试探粒子 $b_*<r<R_*$ 的贡献，并以均匀背景扣除平均势。
令 $\kappa_C=e^2/(4\pi\varepsilon_0)$（单位为 $\si{\joule\metre}$）。
将空间分成小体积元，各元的粒子数独立，均值与方差都为 $n_*\dd V$。
因此未扣平均值之前的势能均值和扣平均之后的方差为
\begin{align}
\overline U&=n_*\int_{b_*}^{R_*}4\pi r^2\frac{\kappa_C}{r}\,\dd r
=2\pi n_*\kappa_C(R_*^2-b_*^2),\label{eq:feqo-mean-coulomb}\\
\sigma_U^2&=n_*\int_{b_*}^{R_*}4\pi r^2
\left(\frac{\kappa_C}{r}\right)^2\dd r
=4\pi n_*\kappa_C^2(R_*-b_*).
\label{eq:feqo-coulomb-energy-spread}
\end{align}
方差的单位是 $\si{\joule\squared}$。
仅写 $n_*[\kappa_C/b_*]^2$ 会留下多余的体积倒数；
碰撞模型若要使用库仑对数，还须独立给出碰撞截面、相对速度和作用时间。

例如 $n_*=10^{18}\,\si{\per\metre\cubed}$、$R_*=\SI{1}{\micro\metre}$、
$b_*=\SI{10}{\nano\metre}$ 给 $\sigma_U=\SI{5.08}{\milli\electronvolt}$。
这只是此静态位置模型的势能涨落，并非某种电镜束流普适的出射能散。
若涨落在时间内保持准静态，对两个耦合系数相差 $\Delta a$ 的通道，
其高斯模型才给 $\exp[-(\Delta a)^2\sigma_U^2t^2/(2\hbar^2)]$。
对一般运动束团，应先计算 $C_U(\tau)=\langle\delta U(\tau)\delta U(0)\rangle$，
再使用已经建立的高斯相位双时间积分；量子关联或非高斯位置统计则需要相应高阶连通函数。
\end{exbox}
